\documentclass{article}
\usepackage[utf8]{inputenc}
\usepackage{amsmath}
\usepackage{geometry}
\geometry{margin=2cm}
\begin{document}

\section*{Resolução da Questão 125 – ENEM 2024 – Ciências da Natureza}

\subsection*{Interpretação do Enunciado e Estratégia}

O enunciado apresenta um procedimento prático adotado por uma agricultora para conservar o sabor adocicado do milho recém-colhido. O ponto central é entender o motivo bioquímico da preservação do sabor com o tratamento térmico descrito: imersão em água fervente seguida de água gelada.

\textbf{Termos-chave:}

\begin{description}
  \item[Desnaturação:] alteração estrutural de enzimas que as tornam inativas, geralmente provocada por calor.
  \item[Enzima:] proteína catalisadora de reações químicas específicas, como a conversão da glicose em amido.
  \item[Choque térmico:] mudança brusca de temperatura, usada para inativar enzimas rapidamente.
\end{description}

O que a questão pede é que se explique por que esse procedimento mantém o sabor adocicado do milho.

\textbf{Estratégia:} entender o papel das enzimas na conversão da glicose em amido e como o calor e o choque térmico interferem nessa reação para preservar a glicose, responsável pelo sabor doce.

\subsection*{Conteúdo, Revisão e Aplicação}

\begin{center}
\begin{tabular}{|l|p{6cm}|p{6cm}|}
\hline
\textbf{Conceito} & \textbf{Ideia-chave} & \textbf{Aplicação no problema} \\
\hline
Enzima & Proteína catalisadora de reações químicas & Enzimas transformam glicose em amido no milho \\
\hline
Desnaturação enzimática & Perda da atividade da enzima por alteração estrutural devida ao calor & Água fervente desnatura enzimas que convertem glicose, bloqueando reação \\
\hline
Choque térmico & Mudança rápida de temperatura que inativa enzimas sem destruir nutrientes & Água gelada após fervura mantém enzimas inativas e integridade do milho \\
\hline
\end{tabular}
\end{center}

\textbf{Fundamentos teóricos:}

As enzimas são essenciais para acelerar reações metabólicas como a conversão de glicose em amido ou outras substâncias. Ao aplicar calor intenso, ocorre a desnaturação, ou seja, as enzimas perdem sua estrutura tridimensional e tornam-se inativas. O choque térmico subsequente com água fria estabiliza essa inativação e preserva o alimento.

\subsection*{Resolução e "Pulo do Gato"}

O procedimento descrito consiste em:

\begin{enumerate}
  \item Imersão em água fervente: a alta temperatura desnatura as enzimas responsáveis por converter glicose em amido.
  \item Imersão em água gelada: evita danos adicionais às células do milho e mantém a inativação enzimática.
\end{enumerate}

\noindent Consequentemente, a glicose não é mais convertida em amido, e o sabor adocicado do milho permanece.

\textbf{Pulo do gato:} o calor da água fervente desnatura as enzimas, bloqueando a conversão da glicose em amido e mantendo o sabor doce – é fundamental focar no efeito das enzimas!

\subsection*{Armadilhas Comuns}

\begin{itemize}
  \item Confundir o papel da temperatura fria e quente.
  \item Pensar que o sabor é mantido por degradação dos nutrientes ou oxidação evitada.
  \item Associar o efeito à desidratação ou remoção de oxigênio.
  \item Ignorar o papel das enzimas no processo.
\end{itemize}

\subsection*{Código da Questão}

CN\_BIO\_ENZ\_001

\vspace{0.5cm}

\noindent Esta resolução segue a estrutura e o conteúdo do gabarito oficial e da análise técnica fornecidos.

\end{document}